\documentclass[11pt]{article}
\usepackage[margin=1in]{geometry}
\usepackage{amsmath,amssymb,amsfonts}
\usepackage{bm}
\usepackage{hyperref}
\usepackage{booktabs}

\DeclareMathOperator{\diag}{diag}
\DeclareMathOperator{\tr}{tr}
\DeclareMathOperator{\softmax}{softmax}
\DeclareMathOperator{\vect}{vec}
\newcommand{\R}{\mathbb{R}}
\newcommand{\E}{\mathbb{E}}
\newcommand{\kron}{\otimes}

\title{Curvature Approximations for Influence Functions:\\A Tutorial on Hessian, GGN, K-FAC, and EK-FAC}
\author{}
\date{}

\begin{document}
\maketitle

\section{Setup}

Consider a classification network $f(x;\theta)\in\R^C$ with parameters
$\theta\in\R^D$ trained by minimising the cross-entropy loss over $N$ samples:
\[
  J(\theta) = \frac{1}{N}\sum_{i=1}^{N} L(z_i,\theta),
  \qquad
  L(z_i,\theta) = -\log \softmax\!\bigl(f(x_i;\theta)\bigr)_{y_i}.
\]
Influence functions require solving systems of the form $\mathbf{H}^{-1}v$
where $\mathbf{H}=\nabla^2_\theta J(\theta^\star)$ is the Hessian at the
trained parameters. Since $\mathbf{H}$ is $D\times D$ and typically
ill-conditioned, a chain of progressively coarser approximations is used.

We study five curvature matrices, each approximating $\mathbf{H}$ with
different fidelity.

\section{The Hessian}

The full Hessian of the empirical risk is
\[
  \mathbf{H} = \frac{1}{N}\sum_{i=1}^{N}\nabla^2_\theta L(z_i,\theta).
\]
We compute it column-by-column using Hessian--vector products (HVPs).
A single HVP $\mathbf{H}v$ is obtained by forward-over-reverse
autodiff: apply \texttt{jax.jvp} to the gradient function.
We batch $K$ HVPs at a time with standard basis vectors to fill $K$
rows of~$\mathbf{H}$, iterating in chunks to control memory.

\section{The Generalised Gauss--Newton (GGN)}

Write the network output as $u_i = f(x_i;\theta)$ with Jacobian
$\mathbf{J}_i = \partial u_i/\partial\theta \in \R^{C\times D}$.
The Hessian decomposes as $\mathbf{H} = \mathbf{G} + \mathbf{R}$
where
\[
  \mathbf{G}
  = \frac{1}{N}\sum_{i=1}^{N}
    \mathbf{J}_i^\top \mathbf{H}_i^{(u)} \mathbf{J}_i,
  \qquad
  \mathbf{H}_i^{(u)}
  = \nabla_u^2\,\phi(u_i,y_i)
  = \diag(p_i) - p_i p_i^\top,
\]
$p_i = \softmax(u_i)$, and the residual $\mathbf{R}$ collects
second-order parameter non-linearities.
The GGN is always positive semi-definite (unlike~$\mathbf{H}$) and
coincides with~$\mathbf{H}$ near convergence where $\mathbf{R}\to 0$.

We compute~$\mathbf{G}$ by forming the per-sample Jacobian
$\mathbf{J}_i\in\R^{C\times D}$ via \texttt{jax.jacobian}, then
accumulating $\mathbf{J}_i^\top \mathbf{H}_i^{(u)} \mathbf{J}_i$
over data chunks.

\section{Block-diagonal GGN (B-GGN)}

Partition parameters by layer:
$\theta=(\theta_1,\dots,\theta_L)$.  The block-diagonal GGN retains
only the diagonal blocks of~$\mathbf{G}$:
\[
  \mathbf{G}_{\text{block}} = \diag(\mathbf{G}_1,\dots,\mathbf{G}_L),
  \qquad
  \mathbf{G}_\ell = \mathbf{G}[\theta_\ell,\theta_\ell].
\]
We compute the full GGN and zero out off-diagonal blocks. The
discarded cross-block terms $\mathbf{G}_{\ell m}$ ($\ell\neq m$)
quantify how parameter groups jointly affect the output.

\section{K-FAC}

\subsection{Per-layer structure}

For a linear layer~$\ell$ without bias, the pre-activation is
$s_\ell = W_\ell\, a_{\ell-1}$ where $a_{\ell-1}\in\R^{d_{\text{in}}}$
is the layer input and $W_\ell\in\R^{d_{\text{out}}\times d_{\text{in}}}$
is the weight.  The per-sample gradient of the loss with respect to the
vectorised weight is
\[
  g_\ell^{(i)} = \vect\!\Bigl(\frac{\partial L_i}{\partial W_\ell}\Bigr)
  = \frac{\partial L_i}{\partial s_\ell} \kron a_{\ell-1}^{(i)}
  = \Delta s_\ell^{(i)} \kron a_{\ell-1}^{(i)},
\]
where we use the row-major (C-order) vectorisation convention matching
JAX's \texttt{ravel}, and $\Delta s_\ell^{(i)} =
\partial L_i/\partial s_\ell$ is the pre-activation gradient.

\medskip
\noindent\textbf{Remark on Kronecker ordering.}\quad
For weights shaped $(d_{\text{out}}, d_{\text{in}})$ with C-order
flattening, the vectorised gradient is
$\Delta s \kron a$ (gradient factor first).
The PyTorch/K-FAC literature uses $(d_{\text{in}}, d_{\text{out}})$
weights with column-major flattening, giving $a\kron\Delta s$.
The resulting Kronecker product for the curvature block is
$\mathbf{S}\kron\mathbf{A}$ (ours) vs.\
$\mathbf{A}\kron\mathbf{S}$ (literature).
Both are correct for their respective conventions.

\subsection{The Kronecker factorisation}

K-FAC approximates each layer's curvature block by assuming
that the input activations $a_{\ell-1}$ and the pre-activation
gradients $\Delta s_\ell$ are statistically independent across
samples:
\[
  \hat{\mathbf{G}}_\ell
  = \mathbf{S}_\ell \kron \mathbf{A}_{\ell-1},
  \qquad
  \mathbf{A}_{\ell-1} = \frac{1}{N}\sum_i a_i a_i^\top,
  \quad
  \mathbf{S}_\ell = \text{(depends on which Fisher)}.
\]

\subsection{Three Fishers and why they matter}
\label{sec:three-fishers}

The ``gradient covariance'' $\mathbf{S}_\ell$ can be defined in three
different ways, leading to different curvature matrices:

\begin{enumerate}
\item \textbf{Empirical Fisher.}\;
  Use the actual data labels~$y_i$:
  \[
    \mathbf{S}_\ell^{\text{emp}}
    = \frac{1}{N}\sum_{i=1}^{N}
      \Delta s_\ell^{(i)}\,\Delta s_\ell^{(i)\top},
    \qquad
    \Delta s_\ell^{(i)}
    = \frac{\partial L(x_i,y_i)}{\partial s_\ell}.
  \]
  Each sample contributes a rank-1 outer product to~$\mathbf{S}_\ell$.

\item \textbf{True (real) Fisher.}\;
  Replace data labels by $\tilde y\sim p(\cdot|x;\theta)$
  sampled from the model's predictive distribution:
  \[
    \mathbf{S}_\ell^{\text{true}}
    = \frac{1}{N}\sum_{i=1}^{N}
      \E_{\tilde y\sim p(\cdot|x_i)}
      \bigl[\Delta s_\ell^{(i,\tilde y)}\,
             \Delta s_\ell^{(i,\tilde y)\top}\bigr]
    = \frac{1}{N}\sum_{i=1}^{N}\sum_{c=1}^{C}
      p_{ic}\;\Delta s_\ell^{(i,c)}\,\Delta s_\ell^{(i,c)\top},
  \]
  where $\Delta s_\ell^{(i,c)}$ is the pre-activation gradient when the
  label is set to class~$c$, and $p_{ic} = p(y{=}c|x_i;\theta)$.

  Each sample now contributes rank up to $C{-}1$ (weighted sum
  over~$C$ classes).

\item \textbf{GGN-derived.}\;
  Define $\mathbf{S}_\ell$ from the GGN block using the sub-network
  Jacobian $\mathbf{B}_\ell = \partial u/\partial s_\ell$:
  \[
    \mathbf{S}_\ell^{\text{GGN}}
    = \frac{1}{N}\sum_{i=1}^{N}
      \mathbf{B}_{\ell,i}^\top\,\mathbf{H}_i^{(u)}\,\mathbf{B}_{\ell,i}.
  \]
\end{enumerate}

\medskip
\noindent
\textbf{Key identity.}\;
For cross-entropy with softmax,
$\mathbf{S}_\ell^{\text{true}} = \mathbf{S}_\ell^{\text{GGN}}$,
because
\[
  \E_{\tilde y\sim p}\bigl[
    (p-e_{\tilde y})(p-e_{\tilde y})^\top
  \bigr]
  = \diag(p) - pp^\top
  = \mathbf{H}^{(u)}.
\]
This identity does \emph{not} hold for the empirical Fisher, which uses a
single label~$y_i$ instead of the expectation.

\medskip
\noindent
\textbf{Practical consequence.}\;
At convergence, the model's predictions become confident
($p\approx e_{y_i}$), so the empirical gradient
$\Delta s \approx 0$ and $\mathbf{S}^{\text{emp}}\to 0$.  The true
Fisher / GGN factor $\mathbf{S}^{\text{true}}$ remains
well-conditioned because the output Hessian
$\diag(p)-pp^\top$ retains curvature across all $C{-}1$
non-predicted classes.

In our experiments, using the empirical Fisher caused K-FAC
approximation errors to explode from ${\sim}10^3$ to ${\sim}10^7$
at 1000 epochs of training.

\subsection{Implementation: analytic vs.\ Monte Carlo}

There are two ways to compute $\mathbf{S}_\ell^{\text{true}}$:

\begin{description}
\item[Analytic (our approach).]
  Compute the sub-network Jacobian
  $\mathbf{B}_\ell = \partial u/\partial s_\ell$ via
  \texttt{jax.jacobian}, then form
  $\mathbf{S}_\ell = \frac{1}{N}\sum_i
   \mathbf{B}_{\ell,i}^\top\,\mathbf{H}_i^{(u)}\,\mathbf{B}_{\ell,i}$.
  This is exact (no sampling noise) and natural for small~$C$.
  The cost is one Jacobian computation per layer.

\item[Monte Carlo (reference implementation).]
  Sample pseudo-labels $\tilde y_i \sim \text{Cat}(p_i)$, run standard
  backprop with these labels, and collect the per-layer gradients
  $\Delta s_\ell^{(i,\tilde y_i)}$ via a \texttt{custom\_vjp} hook.
  Repeat with $M$ Monte Carlo samples and average.  This is unbiased but
  noisy; it requires $M$ backward passes but works for any loss function
  without analytically computing $\mathbf{H}^{(u)}$.
\end{description}

\section{EK-FAC}

K-FAC uses eigenvalues $\Lambda_{\mathbf{S}}\kron\Lambda_{\mathbf{A}}$
(products of marginal spectra) which can misrepresent the true curvature.
EK-FAC keeps the same Kronecker eigenbasis
$\mathbf{U} = \mathbf{U}_{\mathbf{S}}\kron\mathbf{U}_{\mathbf{A}}$ but
replaces the eigenvalues with corrected values that match the diagonal
of the true Fisher / GGN in the Kronecker basis:
\[
  s^\star_{jm}
  = \frac{1}{N}\sum_{i=1}^{N}\sum_{c=1}^{C}
    p_{ic}\;
    \bigl(\mathbf{U}_{\mathbf{S}}^\top
          \Delta s_\ell^{(i,c)}\bigr)_j^2
    \;\bigl(\mathbf{U}_{\mathbf{A}}^\top a_i\bigr)_m^2.
\]
The corrected matrix is
$\hat{\mathbf{G}}_\ell^{\text{EK}} =
 \mathbf{U}\,\diag(\vect(\mathbf{S}^\star))\,\mathbf{U}^\top$.

Since $\Delta s_\ell^{(i,c)} = \mathbf{B}_{\ell,i}^\top(p_i - e_c)$,
we compute the per-class pre-activation gradients from the same
sub-network Jacobian $\mathbf{B}_\ell$ used in K-FAC.

\section{Pseudo-inverse}

All matrices are inverted via eigendecomposition-based pseudo-inverse.
For a symmetric $\mathbf{M} = \mathbf{Q}\Lambda\mathbf{Q}^\top$:
\[
  \mathbf{M}^\dagger
  = \mathbf{Q}\,\Lambda^\dagger\,\mathbf{Q}^\top,
  \qquad
  [\Lambda^\dagger]_{ii}
  = \begin{cases}
      \lambda_i^{-1} & \text{if } |\lambda_i|>\varepsilon,\\
      0 & \text{otherwise},
    \end{cases}
\]
with $\varepsilon=10^{-4}$.  The eigendecomposition is performed in
\texttt{float64} (via NumPy) for numerical stability, then cast
back to \texttt{float32}.

\section{Approximation error metric}

Given test vectors $\{v_i\}$ (per-sample loss gradients), the
approximation error of a curvature estimate~$\hat{\mathbf{H}}$ is:
\[
  \text{Error}(\hat{\mathbf{H}})
  = \frac{1}{N}\sum_{i=1}^{N}
    \frac{\|\mathbf{H}\,\hat{\mathbf{H}}^{\dagger}\,v_i - v_i\|^2}
         {\|v_i\|^2}.
\]
If $\hat{\mathbf{H}} = \mathbf{H}$, the error is zero (up to
numerical precision).  The metric measures how well solving
$\hat{\mathbf{H}}x = v$ approximates solving $\mathbf{H}x = v$.

\section{Summary of the approximation hierarchy}

\begin{center}
\begin{tabular}{lll}
  \toprule
  \textbf{Method} & \textbf{What it drops} & \textbf{Our computation} \\
  \midrule
  Hessian   & (nothing)               & HVPs via \texttt{jax.jvp} of \texttt{jax.grad} \\
  GGN       & Residual $\mathbf{R}$   & $\mathbf{J}^\top\mathbf{H}^{(u)}\mathbf{J}$ via \texttt{jax.jacobian} \\
  B-GGN     & Cross-layer blocks      & Diagonal blocks of GGN \\
  EK-FAC    & Off-diagonal in Kron.\ basis & Corrected eigenvalues in $\mathbf{U}_S\kron\mathbf{U}_A$ \\
  K-FAC     & Eigenvalue correction   & $\mathbf{S}\kron\mathbf{A}$ with marginal spectra \\
  \bottomrule
\end{tabular}
\end{center}

Each step discards information: the residual (non-linearity), cross-layer
coupling, and within-block Kronecker structure.  The paper's central
finding is that better curvature fidelity consistently improves influence
function quality.

\end{document}
